% !TEX spellcheck = en_US
% !TEX spellcheck = LaTeX
\documentclass[letterpaper,10pt,english]{article}
\usepackage{%
amsfonts,%
amsmath,%
amssymb,%
amsthm,%
babel,%
bbm,%
%biblatex,%
caption,%
centernot,%
color,%
enumerate,%
%enumitem,%
epsfig,%
epstopdf,%
etex,%
fancybox,%
framed,%
fullpage,%
%geometry,%
graphicx,%
hyperref,%
latexsym,%
mathptmx,%
mathtools,%
multicol,%
paralist,%
pgf,%
pgfplots,%
pgfplotstable,%
pgfpages,%
proof,%
psfrag,%
%subfigure,%
tcolorbox,%
tfrupee,%
tikz,%
times,%
%ulem,%
url,%
xcolor,%
mathpazo,%
}

\definecolor{shadecolor}{gray}{.95}%{rgb}{1,0,0}
\usepackage[margin=1in,top=0.75in]{geometry}
\usepackage[mathscr]{eucal}
\usepgflibrary{shapes}
\usepgfplotslibrary{fillbetween}
\usetikzlibrary{%
arrows,%
backgrounds,%
chains,%
decorations.pathmorphing,% /pgf/decoration/random steps | erste Graphik
decorations.text,% 
matrix,%
positioning,% wg. " of "
fit,%
patterns,%
petri,%
plotmarks,%
scopes,%
shadows,%
shapes.misc,% wg. rounded rectangle
shapes.arrows,%
shapes.callouts,%
shapes%
}

%\pgfplotsset{compat=newest} %<------ Here
\pgfplotsset{compat=1.11} %<------ Or use this one

\theoremstyle{plain}
\newtheorem{thm}{Theorem}[section]
\newtheorem{lem}[thm]{Lemma}
\newtheorem{prop}[thm]{Proposition}
\newtheorem{cor}[thm]{Corollary}
\newtheorem{clm}[thm]{Claim}

\theoremstyle{definition}
\newtheorem{defn}[thm]{Definition}
\newtheorem{conj}[thm]{Conjecture}
\newtheorem{exmp}[thm]{Example}
\newtheorem{axiom}[thm]{Axiom}
\newtheorem{assum}[thm]{Assumption}
\newtheorem{exerc}[thm]{Exercise}
\newtheorem*{defin}{Definition}

\theoremstyle{remark}
\newtheorem{rem}{Remark}
\newtheorem{note}{Note}

\newcommand{\iid}{\emph{i.i.d.}\ }
\newcommand{\Cov}{\operatorname{Cov}}
%\newcommand{\det}{\operatorname{det}}
%\newcommand{\Real}{\mathbb{R}}
\newcommand{\tr}{\operatorname{tr}}
\newcommand{\Var}{\operatorname{Var}}
\newcommand{\cov}{\operatorname{cov}}

%\renewcommand{\proof}[1]{\begin{proof}#1\end{proof}}
\newcommand{\EQ}[1]{\begin{equation*}#1\end{equation*}}
\newcommand{\EQN}[1]{\begin{equation}#1\end{equation}}
\newcommand{\eq}[1]{\begin{align*}#1\end{align*}}
\newcommand{\meq}[2]{\begin{xalignat*}{#1}#2\end{xalignat*}}
\newcommand{\norm}[1]{\left\lVert#1\right\rVert}
\newcommand{\inner}[1]{\left\langle #1\right\rangle}
\newcommand{\abs}[1]{\left\lvert#1\right\rvert}
\newcommand{\expect}[1]{\mathbb{E}\left[{#1}\right]}
\newcommand{\prob}[1]{\mathbb{P}\left[{#1}\right]}
\newcommand{\given}{\; \big\vert \;} 
\newcommand{\set}[1]{\left\{#1\right\}} 
\newcommand{\indicator}[1]{1_{\set{#1}}} 
\newcommand{\Ind}[1]{\mathbbm{1}_{#1}} 
\newcommand{\SetIn}[1]{\mathbbm{1}_{\set{#1}}} 
\newcommand{\red}[1]{\textcolor{red}{#1}} 
\newcommand{\floor}[1]{\left\lfloor#1\right\rfloor}
\newcommand{\ceil}[1]{\left\lceil#1\right\rceil}


\newcommand{\C}{\mathbb{C}}
\newcommand{\D}{\mathbb{D}}
\newcommand{\E}{\mathbb{E}}
\newcommand{\F}{\mathbb{F}}
\newcommand{\N}{\mathbb{N}}
\renewcommand{\P}{\mathbb{P}}
\newcommand{\Q}{\mathbb{Q}}
\newcommand{\R}{\mathbb{R}}
\newcommand{\Z}{\mathbb{Z}}

\newcommand{\bA}{\mathbf{A}}
\newcommand{\bI}{\mathbf{I}}
\newcommand{\bU}{\mathbf{1}}
\newcommand{\bx}{\mathbf{x}}
\newcommand{\bX}{\mathbf{X}}
\newcommand{\bZ}{\mathbf{Z}}


\newcommand{\cA}{\mathcal{A}}
\newcommand{\cB}{\mathcal{B}}
\newcommand{\cC}{\mathcal{C}}
\newcommand{\cF}{\mathcal{F}}
\newcommand{\cG}{\mathcal{G}}
\newcommand{\cM}{\mathcal{M}}
\newcommand{\cN}{\mathcal{N}}
\newcommand{\cP}{\mathcal{P}}
\newcommand{\cT}{\mathcal{T}}
\newcommand{\cX}{\mathcal{X}}
\newcommand{\cY}{\mathcal{Y}}

\newcommand{\sA}{\mathscr{A}}
\newcommand{\sB}{\mathscr{B}}
\newcommand{\sC}{\mathscr{C}}
\newcommand{\sE}{\mathscr{E}}
\newcommand{\sF}{\mathscr{F}}
\newcommand{\sG}{\mathscr{G}}
\newcommand{\sH}{\mathscr{H}}
\newcommand{\sL}{\mathscr{L}}
\newcommand{\sS}{\mathscr{S}}
\newcommand{\sT}{\mathscr{T}}
\newcommand{\sX}{\mathscr{X}}
\newcommand{\sY}{\mathscr{Y}}
\newcommand{\sZ}{\mathscr{Z}}

\newcommand{\tS}{\tilde{S}}
% Debug
\newcommand{\todo}[1]{\begin{color}{blue}{{\bf~[TODO:~#1]}}\end{color}}

% a few handy macros

\renewcommand{\le}{\leqslant}
\renewcommand{\ge}{\geqslant}
\newcommand\matlab{{\sc matlab}}
\newcommand{\goto}{\rightarrow}
\newcommand{\bigo}{{\mathcal O}}
%\newcommand{\half}{\frac{1}{2}}
%\newcommand\implies{\quad\Longrightarrow\quad}
\newcommand\reals{{{\rm l} \kern -.15em {\rm R} }}
\newcommand\complex{{\raisebox{.043ex}{\rule{0.07em}{1.56ex}} \hskip -.35em {\rm C}}}


% macros for matrices/vectors:

% matrix environment for vectors or matrices where elements are centered
\newenvironment{mat}{\left[\begin{array}{ccccccccccccccc}}{\end{array}\right]}
\newcommand\bcm{\begin{mat}}
\newcommand\ecm{\end{mat}}

% matrix environment for vectors or matrices where elements are right justifvied
\newenvironment{rmat}{\left[\begin{array}{rrrrrrrrrrrrr}}{\end{array}\right]}
\newcommand\brm{\begin{rmat}}
\newcommand\erm{\end{rmat}}

% for left brace and a set of choices
%\newenvironment{choices}{\left\{ \begin{array}{ll}}{\end{array}\right.}
\newcommand\when{&\text{if~}}
\newcommand\otherwise{&\text{otherwise}}
% sample usage:
%\delta_{ij} = \begin{choices} 1 \when i=j, \\ 0 \otherwise \end{choices}


% for labeling and referencing equations:
\newcommand{\eql}{\begin{equation}\label}
\newcommand{\eqn}[1]{(\ref{#1})}
% can then do
%\eql{eqnlabel}
%...
%\end{equation}
% and refer to it as equation \eqn{eqnlabel}.
% some useful macros for finite difference methods:
\newcommand\unp{U^{n+1}}
\newcommand\unm{U^{n-1}}
% for chemical reactions:
\newcommand{\react}[1]{\stackrel{K_{#1}}{\rightarrow}}
\newcommand{\reactb}[2]{\stackrel{K_{#1}}{~\stackrel{\rightleftharpoons}
 {\scriptstyle K_{#2}}}~}
\makeatletter
\def\th@plain{%
\thm@notefont{}% same as heading font
\itshape % body font
}
\def\th@definition{%
\thm@notefont{}% same as heading font
\normalfont % body font
}
\makeatother
\date{}


\title{Lecture-16: Weak convergence}% of random variables}
\author{}

\begin{document}
\maketitle

%Consider a random sequence $X:\Omega\to\R^\N$ defined on a probability space $(\Omega, \sF, P)$.
\section{Convergence in distribution}
\begin{defn}[convergence in distribution] 
A  random sequence $X:\Omega\to\R^\N$ defined on a probability space $(\Omega, \sF, P)$ \textbf{converges in distribution} to a random variable $X_\infty:\Omega'\to\R$ defined on a probability space $(\Omega', \sF', P')$  if $\lim_nF_{X_n}(x) = F_{X_\infty}(x)$ at all continuity points $x$ of $F_{X_\infty}$. 
Convergence in distribution is denoted by $\lim_nX_n = X_\infty$ in distribution. 
\end{defn}
\begin{prop}
\label{Prop:ConvDistEquivalence}
Consider a random sequence$X:\Omega\to\R^\N$ defined on a probability space $(\Omega, \sF, P)$ and a random variable $X_\infty: \Omega'\to\R$ defined on another probability space $(\Omega', \sF', P')$.  
%, both defined on the same . 
Then the following statements are equivalent.
\begin{enumerate}[(a)]
\item $\lim_nX_n = X_\infty$ in distribution.
\item $\lim_n\E[g(X_n)] = \E[g(X_\infty)]$ for any bounded continuous Borel measurable function $g:\R\to\R$.
\item Characteristic functions converge point-wise, i.e. 
$\lim_n\Phi_{X_n}(u) = \Phi_{X_\infty}(u)$ for each $u \in \R$. 
\end{enumerate}
\end{prop}
\begin{proof}
Let $X:\Omega\to\R^\N$ be a sequence of random variables and let $X_\infty:\Omega'\to\R$ be a random variable. 
We will show that $(a) \implies (b) \implies (c) \implies (a)$.
\begin{enumerate}[(a)]
\item [$(a) \implies (b)$:] 
%Let $\lim_nX_n = X_\infty$ in distribution, then 
Applying the bounded convergence theorem to any bounded continuous Borel measurable function $g:\R\to\R$, we have $\lim_n\int_{x \in \R} g(x)dF_{X_n}(x) = \int_{x \in \R} g(x)\lim_ndF_{X_n}(x)$.
\item [$(b) \implies (c)$:] 
%Let $\lim_n\E[g(X_n)] = \E[g(X_\infty)]$ for any bounded continuous function $g$. 
Taking $g(x) = e^{jux}$, we get the result. 
\item [$ (c) \implies (a)$:] 
The proof of this part is technical and is omitted.
%Let $\lim_n\Phi_{X_n}(u) \to \Phi_X(u)$ for each $u \in \R$. 
\end{enumerate}
\end{proof}
\begin{shaded}
\begin{exmp}[Convergence in distribution but not in probability] 
Consider a sequence of non-degenerate continuous \iid random variables $X:\Omega\to\R^\N$ and independent random variable $Y: \Omega'\to\R$, all with the common distribution $F_Y$. 
Then $F_{X_n} = F_Y$ for all $n \in \N$, and hence $\lim_nX_n = Y$ in distribution. 
If the common distribution $F_Y$ is zero mean Gaussian with variance $\sigma^2$, 
then $X_n - Y$ is zero mean Gaussian with variance $2\sigma^2$. 
Therefore, for $\epsilon < \sigma\sqrt{\pi}$ and all $n \in \N$
\EQ{
P\set{\abs{X_n - Y} \le \epsilon} 
= \frac{1}{\sqrt{4\pi\sigma^2}}\int_{z \in [-\epsilon, \epsilon]}e^{-z^2/4\sigma^2}dz
\le \frac{\epsilon}{\sigma\sqrt{\pi}} < 1.
}
It follows that $P\set{\abs{X_n - Y} > \epsilon} \ge 1- \frac{\epsilon}{\sigma\sqrt{\pi}}$ for all $n \in \N$, and hence $\lim X_n \neq Y$ in probability.
%However, for any $n \in \N$ and $\epsilon > 0$, from the monotonicity of distribution function, we have 
%\EQ{
%P\set{\abs{X_n - Y} \le \epsilon} 
%= \E\SetIn{X_n \in [Y-\epsilon, Y+\epsilon]} 
%= \E F_Y(Y+\epsilon) - \E F_Y(Y-\epsilon) < 1.
%}
%Defining $Z_1 \triangleq F_Y(Y+\epsilon)$ and $Z_2 \triangleq F_Y(Y-\epsilon)$, 
%we observe that 
%\begin{xalignat*}{2}
%&F_{Z_1}(z) %= P\set{Y \le F_Y^{-1}(z)-\epsilon} 
%= F_Y(F_Y^{-1}(z)-\epsilon)
%= z - \epsilon F'_Y(F_Y^{-1}(z)),&
%&F_{Z_2}(z) %= P\set{Y \le F_Y^{-1}(z)+\epsilon} 
%= F_Y(F_Y^{-1}(z)+\epsilon)
%= z + \epsilon F'_Y(F_Y^{-1}(z).
%\end{xalignat*}
\end{exmp}
\end{shaded}
\begin{lem}[Convergence in probability implies in distribution] 
Consider a sequence $X:\Omega\to\R^\N$ of random variables and a random variable $X_\infty: \Omega\to\R$ defined on a probability space $(\Omega, \sF, P)$, 
such that $\lim_nX_n = X_\infty$ in probability, then $\lim_nX_n = X_\infty$ in distribution.
\end{lem}
\begin{proof} 
We will show that all continuity points $x$ of $F_{X_\infty}$, 
we have $\lim_{n\to\infty}F_{X_n}(x) = F_{X_\infty}(x)$. 
Fix $\epsilon > 0$. 
Since $x$ is a continuity point of non-decreasing function $F_{X_\infty}$, 
choose $\delta > 0$ such that $F_{X_\infty}(x+\delta) - F_{X_\infty}(x-\delta) < \epsilon$. 
Therefore, it suffices to show that 
\EQ{
F_{X_\infty}(x-\delta) \le \lim\inf_{n\to\infty}F_{X_n}(x)\le \lim\sup_{n\to\infty}F_{X_n}(x) \le F_{X_\infty}(x+\delta).
}
%Consider the events $A_n(x) \triangleq \set{X_n\le x}$ and $A(x) \triangleq \set{X \le x}$, then we can write  
%\EQ{
%A_n(x) = (A_n(x)\cap A(x+\epsilon)) \cup (A_n(x)\cap A^c(x+\epsilon))
%}
%Fix $\epsilon > 0$,
For the chosen $\delta >0$, we consider the event  $A_n(\delta) \triangleq \set{\omega \in \Omega: \abs{X_n(\omega)-X_\infty(\omega)} > \delta} 
= \set{X_n \notin [X_\infty- \delta, X_\infty+\delta]}\in \sF,$
and define events $A_{X_n}(x) \triangleq \set{X_n \le x}$ and $A_{X_\infty}(x) \triangleq \set{X_\infty \le x}$. 
Then, we can write
\begin{xalignat*}{2}
&A_{X_n}(x)\cap A_{X_\infty}(x+\delta) \subseteq A_{X_\infty}(x+\delta),&&A_{X_n}(x)\cap A_{X_\infty}^c(x+\delta) \subseteq A_n(\delta),\\
&A_{X_\infty}(x-\delta)\cap A_{X_n}(x) \subseteq A_{X_n}(x), &&A_{X_\infty}(x-\delta)\cap A_{X_n}^c(x) \subseteq A_n(\delta).
\end{xalignat*}
From the above set relations, law of total probability, and union bound, we have 
\EQ{
F_{X_\infty}(x-\delta) - P(A_n(\delta)) \le F_{X_n}(x) \le F_{X_\infty}(x+\delta) + P(A_n(\delta)).
}
From the convergence in probability, we have $\lim_nP(A_n(\delta)) = 0$, and the result follows. 
%Therefore, we get 
%\EQ{
%F_{X_\infty}(x-\delta) \le \lim\inf_nF_{X_n}(x) \le  \lim\sup_nF_{X_n}(x) \le F_{X_\infty}(x+\delta). 
%}
%We get the result at the continuity points of $F_{X_\infty}$, since the choice of $\epsilon$ was arbitrary. 
\end{proof}
\begin{thm}[Central Limit Theorem] 
Consider an \iid random sequence $X:\Omega\to\R^\N$ defined on a probability space $(\Omega, \sF, P)$, 
with $\E X_n = \mu$ and $\Var(X_n) = \sigma^2$ for all $n\in\N$.  
We define the $n$-sum as $S_n\triangleq \sum_{i=1}^nX_i$ and consider a standard normal random variable $Y:\Omega\to\R$ with density function $f_Y(y)= \frac{1}{\sqrt{2\pi}}e^{-\frac{y^2}{2}}$ for all $y \in \R$. 
Then, 
\EQ{
\lim_n\frac{S_n-n\mu}{\sigma\sqrt{n}} = Y\text{ in distribution.}
}
\end{thm}
\begin{proof} 
The classical proof is using the characteristic functions. 
Let $Z_i \triangleq \frac{X_i-\mu}{\sigma}$ for all $i \in \N$, 
then the shifted and scaled $n$-sum is given by $\frac{S_n-n\mu}{\sigma\sqrt{n}} = \frac{1}{\sqrt{n}}\sum_{i=1}^nZ_i$. 
We use the third equivalence in Proposition~\ref{Prop:ConvDistEquivalence} to show that the characteristic function of converges to the characteristic function of the standard normal. 
We define the characteristic functions 
\begin{xalignat*}{3}
&\Phi_{n}(u) \triangleq \E\exp\left(ju\frac{(S_n-n\mu)}{\sigma\sqrt{n}}\right), &
&\Phi_{Z_i}(u) \triangleq \E\exp(juZ_i), &
&\Phi_Y(u) \triangleq \E\exp(juY).
\end{xalignat*}
We can compute the characteristic function of the standard normal as 
\EQ{
\Phi_Y(u) = \frac{1}{\sqrt{2\pi}}\int_{y \in \R}e^{-\frac{u^2}{2}}\exp\left(-\frac{(y-ju)^2}{2}\right)dy = e^{-\frac{u^2}{2}}.
}
Since the random sequence $Z: \Omega\to \R^\N$ is a zero mean \iid sequence, 
it follows that $\Phi_{Z_1}^{(1)}(0) = j\E Z_1 = 0$ and $\Phi_{Z_1}^{(2)}(0) = j^2\E Z_1^2 = - 1$.  
Using the Taylor expansion of the characteristic function $\Phi_{Z_1}$, we have 
\EQ{
\Phi_{n}(u) = \prod_{i=1}^n\E\exp\left(ju\frac{(X_i-\mu)}{\sigma\sqrt{n}}\right) = \left[\Phi_{Z_1}\left(\frac{u}{\sqrt{n}}\right)\right]^n = \left[1 - \frac{u^2}{2n} + o\left(\frac{u^2}{n}\right)\right]^n. 
}
For any $u \in \R$, taking limit $n \in \N$, we get the result.
\end{proof}

\section{Strong law of large numbers}
\begin{defn} 
For a random sequence $X:\Omega\to\R^\N$ defined on a probability space $(\Omega,\sF,P)$ with bounded mean $\E\abs{X_n} < \infty$ for all $n \in \N$, 
we define the $n$-sum as $S_n \triangleq \sum_{i=1}^nX_i$ and the empirical $n$-mean $\frac{S_n}{n}$ for each $n \in \N$. 
For each $n \in \N$, we define event 
\EQ{
E_n \triangleq \set{\omega \in \Omega: \abs{S_n - \E S_n} > n\epsilon} \in \sF.
} 
\end{defn}
\begin{thm}[$L^4$ strong law of large numbers] 
Let $X:\Omega\to\R^\N$ be a sequence of independent random variables defined on probability space $(\Omega,\sF,P)$ with bounded mean $\E X_n$ for each $n \in \N$ and uniformly bounded fourth central moment $\sup_{n \in \N}\E(X_n-\E X_n)^4 \le B < \infty$. %for all $n \in \N$.  
Then, the empirical $n$-mean converges to $\lim_n\frac{\E S_n}{n}$ almost surely.
\end{thm}
\begin{proof}
%Applying Jensen's inequality to convex function $g:\R \to \R$ defined by $x \mapsto x^4$, 
%and to random variable $Z$ that takes value $n(X_i-\E X_i)$ with probability $\frac{1}{n}$ for all $i \in [n]$, 
%we get 
%\EQ{
%(S_n-\E S_n)^4 
%\le n^3\sum_{i=1}^n(X_i-\E X_i))^4.
%}
Recall that $\E(S_n-\E S_n)^4 = \E(\sum_{i=1}^n(X_i-\E X_i))^4 = \sum_{i=1}^n\E(X_i-\E X_i)^4 + 3\sum_{i=1}^n\sum_{j \neq i}\E(X_i-\E X_i)^2\E(X_j-\E X_j)^2$. 
Recall that when the fourth moment is bounded, then so is second moment. 
Hence, $\sup_{i\in \N}\E(X_i-\E X_i)^2 \le C$ for some $C \in \R_+$.
Therefore, from the Markov's inequality, we have %and Minkowski's inequality, we have 
\EQ{
P(E_n) 
\le \frac{\E(S_n-\E S_n)^4}{n^4\epsilon^4} 
%\le \frac{\sum_{i=1}^{n}\E(X_i-\mu)^4}{n^4\epsilon^4} 
\le \frac{nB+3n(n-1)C^2}{n^4\epsilon^4}.
} 
It follows that the $\sum_{n \in \N}P(E_n) < \infty$, and hence by Borel Canteli Lemma, we have 
\EQ{
P\set{E_n^c \text{ for all but finitely many }n} = 1.
} 
Since, the choice of $\epsilon$ was arbitrary, the result follows.
\end{proof}
\begin{thm}[$L^2$ strong law of large numbers] 
Let $X:\Omega \to \R^\N$ be a sequence of pair-wise uncorrelated random variables defined on a probability space $(\Omega, \sF, P)$ with bounded mean $\E X_n$ for all $n \in \N$ and uniformly bounded variance $\sup_{n\in\N}\Var(X_n) \le B < \infty$. %for all $n \in \N$.  
Then, the empirical $n$-mean converges to $\lim_n\frac{\E S_n}{n}$ almost surely.
\end{thm}
\begin{proof} 
For each $n \in \N$, we define events $F_n \triangleq E_{n^2}$, and 
\EQ{
G_n \triangleq \set{\max_{n^2 \le k < (n+1)^2}\abs{S_k-S_{n^2} - \E(S_k-S_{n^2})} > n^2\epsilon} 
= \bigcup_{n^2 \le k < (n+1)^2}\set{\omega \in \Omega: \abs{\sum_{i=n^2}^{k}(X_i - \E X_i)} > n^2\epsilon}.
}
From the Markov's inequality and union bound, we have 
\meq{2}{
&P(F_n) 
\le \frac{\sum_{i=1}^{n^2}\Var(X_i)}{n^4\epsilon^2} 
\le \frac{B}{n^2\epsilon^2},&
&P(G_n) 
\le \sum_{k=n^2}^{(n+1)^2-1}\frac{(k-n^2+1)B}{n^4\epsilon^2} 
\le \frac{(2n+1)^2B}{n^4\epsilon^2}.
}
Therefore, $\sum_{n \in \N}P(F_n) < \infty$ and $\sum_{n \in \N}P(G_n) < \infty$, 
and hence by Borel Canteli Lemma, 
we have 
\EQ{
\lim_n\frac{S_{n^2}-\E S_{n^2}}{n^2} = \lim_n\max_{n^2 \le k < (n+1)^2-1}\frac{S_k-S_{n^2}-\E(S_k-S_{n^2})}{n^2} = 0 \text{ a.s. }
}
The result follows from the fact that for any $k \in \N$, there exists $n \in \N$  such that $k \in \set{n^2, \dots, (n+1)^2-1}$ and hence 
\EQ{
\frac{\abs{S_k-\E S_k}}{k} \le \left(
\frac{\abs{S_{n^2}-\E S_{n^2}}}{n^2} + \frac{\abs{S_k-S_{n^2}-\E(S_k-S_{n^2})}}{n^2}
\right).
}
\end{proof}
\begin{thm}[$L^1$ strong law of large numbers] 
Let $X:\Omega\to\R^\N$ be a random sequence defined on a probability space $(\Omega,\sF,P)$ such that $\sup_{n\in\N}\E\abs{X_n} \le B < \infty$. %for all $n \in \N$.  
Then, the empirical $n$-mean converges to $\lim_n\frac{\E S_n}{n}$ almost surely.
\end{thm}
%\begin{proof} 
%\red{
%For each $n \in \N$, we define truncated random variable $Y_n \triangleq X_n\SetIn{\abs{X_n} \le n}$. 
%Then, we can write 
%\EQ{
%\sum_{n \in \N}P\set{X_n \neq Y_n} = \sum_{n\in \N}P\set{\abs{X_n} > n} \le \E\abs{X_1} < \infty.
%}
%We will first show that $\lim_{n}\frac{1}{n}\sum_{i=1}^nY_i = \E X_1$ a.s. 
%Towards this end, we define event 
%\EQ{
%E_n \triangleq \set{\omega \in \Omega: \frac{1}{n}\abs{\sum_{i=1}^n(Y_i(\omega) - \E Y_i)} > \epsilon} \in \sF.
%} 
%From the Markov's inequality, we have $P(E_n) \le \frac{\sum_{i=1}^n\Var(Y_i)}{n^2\epsilon^2}$.
%}
%\end{proof}

\end{document}